\section{Introduction}

A mean-payoff game asks whether a player can maintain a nonnegative
long-run average reward while an opponent chooses some of the moves.
The arena is finite, but plays are infinite and integer edge weights may
have large magnitudes. Positional strategies provide finite certificates
for the two winning regions. The algorithmic difficulty is to find those
regions with a cost controlled by the binary description, rather than by
the numerical size of a weight.

This paper gives a randomized quasipolynomial-time algorithm. Its central
output is an exact bounded integer potential: with probability at least
$7/8$, every coordinate of that potential is correct. Thresholding it
returns the complete winning set. The method also gives order guarantees
in finite integer boxes with arbitrary endpoints and the allowed reference
width. Each guarantee concerns a fixed comparison vector and all its
coordinates at once.

\subsection{Games, encoding, and the main theorem}

The input consists of a finite directed graph $G=(V,E)$ on
$V=\{1,\ldots,n\}$, $n\ge1$, an ownership partition
$V=V_{\mathrm{Max}}\mathbin{\dot\cup}V_{\mathrm{Min}}$, and signed
integer weights $w:E\to\Z$. Every vertex has an outgoing edge;
self-loops and parallel edges are permitted. At a vertex its owner
chooses the next edge. Both players may use strategies depending on the
entire finite history. For a play $\pi=(e_0,e_1,\ldots)$, define
\[
 \operatorname{MP}_w(\pi)
 =\liminf_{T\to\infty}\frac1T\sum_{j=0}^{T-1}w(e_j).
\]
We use an explicit binary encoding of the graph, ownership, and signed
weights. Its total length is $L$; in particular, endpoints and ownership
are charged as part of the input.

\begin{theorem}\label{thm:main}
There is a uniform randomized algorithm that, for every such game,
returns the entire set
\[
 U=\{v\in V:\text{there is a Max strategy }\sigma
       \text{ such that for every Min strategy }\tau,\ 
       \operatorname{MP}_w(\pi_{v,\sigma,\tau})\ge0\}
\]
correctly with probability at least $7/8$.
For an absolute constant $C>0$, on every random tape its total cost is at
most
\[
 2^{C(\log_2(L+2))^2}
\]
bit operations in a standard model polynomially equivalent to a uniform
probabilistic Turing machine. The bound includes preprocessing, recursive
calls, intermediate arithmetic, and random-bit generation.
\end{theorem}

The probability statement concerns the whole returned set, and the time
bound holds even when that output is wrong. Large weights are allowed:
the resource parameter is their binary description together with the
rest of the explicit input. The asserted bound is quasipolynomial;
a polynomial-time bound would be stronger.

The recovered potential has a polynomial-time check. A run can therefore
return either a verified winning set with positional strategies for both
players, or a failure symbol. It returns the verified output with
probability at least $7/8$ and never accepts an incorrect one.
Corollary~\ref{cor:certified-output} proves this consequence and an
always-correct, or Las Vegas, version obtained by repeating until the
check succeeds. The repeated version has quasipolynomial expected bit
cost; each individual trial retains the bound on every tape. The
strategies certify the zero-threshold regions, with Max enforcing
nonnegative payoff on its region and Min enforcing strictly negative
payoff on the complement.

The companion article \emph{Deterministic quasipolynomial-time
mean-payoff games}~\cite[Theorem~1.1]{OpenAI2026Deterministic} gives
a deterministic complete winning-set algorithm with the same form of
bit bound. Its Section~7 also reduces exact values and globally optimal
positional strategies to polynomially many winning-set queries. The
value at a start is the best worst-case mean payoff from that start;
the two strategies attain these values simultaneously at all starts.
Using amplified calls to the algorithm of Theorem~\ref{thm:main},
that reduction returns
the entire value and strategy output correctly with probability at
least $7/8$, with the same form of bit bound on every random tape
\cite[Section~7]{OpenAI2026Deterministic}.

The present paper retains different outputs and a different method:
pointwise whole-vector comparisons in boxes with arbitrary integer
endpoints and allowed reference width, exact recovery of the top-box
potential, and a randomized orientation recursion. The proof of
Theorem~\ref{thm:main} is independent of the deterministic algorithm.

\subsection{Historical context and related algorithms}

Ehrenfeucht and Mycielski proved the positional structure of mean-payoff
games~\cite{EhrenfeuchtMycielski1979}. A long algorithmic line uses
\emph{potentials}, numbers assigned to vertices whose endpoint
differences reweight edges without changing cycle sums. Gurvich,
Karzanov, and Khachiyan developed this approach for cyclic games
\cite{GurvichKarzanovKhachiyan1988}; related algorithms and analyses
were given by Karzanov and Lebedev~\cite{KarzanovLebedev1993} and
Pisaruk~\cite{Pisaruk1999}. Zwick and Paterson established complexity
results and pseudopolynomial algorithms for mean-payoff threshold and
value problems~\cite{ZwickPaterson1996}. A pseudopolynomial cost may be
polynomial in the weight magnitude while exponential in its binary
length, which is the distinction relevant here.

Energy games express a related finite-prefix requirement: starting with
some credit, keep every cumulative total nonnegative. Resource
interfaces and energy constraints provide a verification setting for
this question~\cite{ChakrabartiEtAl2003,BouyerEtAl2008}. Brim,
Chaloupka, Doyen, Gentilini, and Raskin use bounded integer progress
measures and monotone lifting~\cite{BrimEtAl2011}; their threshold
algorithm has $O(|E|nM)$ arithmetic cost for an absolute weight bound
$M$~\cite[Theorem~8]{BrimEtAl2011}. Lifshits and Pavlov give a paired
potential characterization and an exponential algorithm
\cite{LifshitsPavlov2007}. Cadilhac, Casares, and Ohlmann develop a
uniform account of sound potential updates and accelerated value
iteration for energy games~\cite{CadilhacCasaresOhlmann2025}.
Our box comparison method uses this potential viewpoint, with a
symmetric boundary convention and a locally proved game reduction.

Strategy improvement supplies another route. For graphs without
parallel edges, Bj\"orklund and Vorobyov obtain an expected
$2^{O(\sqrt{n\log n})}$ arithmetic-operation bound for the threshold
partition~\cite[Theorem~7.1]{BjorklundVorobyov2007}, following earlier
work with Sandberg~\cite{BjorklundSandbergVorobyov2004}. Andersson and
Vorobyov use discounted linear programs for randomized strongly
subexponential value and strategy algorithms
\cite{AnderssonVorobyov2006}. Subsequent analyses include Zwick's
bounds for Random-Action-Removal~\cite{Zwick2026} and the bounds of
Dorfman, Kaplan, and Zwick for Switch-All and Random-Edge in energy
games~\cite{DorfmanKaplanZwick2026}. These comparisons require two
distinctions. Arithmetic counts must be supplemented by operand bit
lengths. Moreover, an expected Las Vegas bit bound can be converted to
a bounded-error bound on every tape by a fixed cutoff. Thus the
quantitative contribution of Theorem~\ref{thm:main} is its
quasipolynomial dependence on $L$, together with the specified complete
output and potential guarantees.

The 2017 quasipolynomial parity-game algorithm of Calude, Jain,
Khoussainov, Li, and Stephan~\cite{CaludeJainKhoussainovLiStephan2022}
opened a further direction. In a parity game, the winner is determined
by vertex priorities appearing infinitely often. Extensions to
quantitative objectives retain additional parameters: Daviaud,
Jurdzi\'nski, and Lazi\'c combine succinct parity progress measures
with energy bounds in a pseudo-quasipolynomial algorithm for
mean-payoff parity games~\cite{DaviaudJurdzinskiLazic2018}.
Fijalkow, Gawrychowski, and Ohlmann obtain universal-graph bounds
involving numerical weight magnitude or an exponential dependence on
the number of distinct weights~\cite{FijalkowGawrychowskiOhlmann2020}.
The later theory also proves mean-payoff universal-graph size lower
bounds~\cite[Section~7]{ColcombetFijalkowGawrychowskiOhlmann2022}.
Such lower bounds concern that representation of the objective and do
not bound the complexity of every algorithm.

Randomness in the input leads to a different comparison. Loff and
Skomra prove high-probability polynomial bounds for independent Gaussian
perturbations of normalized payoffs on ergodic game graphs
\cite[Theorem~1]{LoffSkomra2024}. Their bound also depends on the
perturbation parameter. Here the input is arbitrary and fixed; only
the algorithm is random.

A dominion in a parity game is a vertex set in which one player can
keep the play and win. Parys's recursion uses separate bounds on the
sizes of the two players' dominions. Among calls with a reduced bound, it makes one call
retaining both current bounds
\cite[Algorithm~2 and Section~5]{Parys2019}. Recursive refinements
are presented by Lehtinen, Parys, Schewe, and Wojtczak
\cite[Sections~3--4]{LehtinenParysScheweWojtczak2022}. Our separate
support budgets parallel these bounds on protected objects. Here the
objects are potential vectors, and each probabilistic comparison is
stated for one fixed vector; the parity-game contracts cover all
qualifying dominions simultaneously.

Ohlmann's symmetric recursion for mean-payoff games supplies a related
potential viewpoint and leaves its subexponential complexity analysis
open \cite[Sections~3--5 and~7]{Ohlmann2026}. We prove the weighted
support invariant, the sampled orientation argument, and the resulting
bit bound below, without importing a recursive solver as an oracle.

Akian, Gaubert, and Guterman establish reductions between tropical
feasibility and winning initial states in mean-payoff games
\cite[Theorem~3.2 and Corollary~3.7]{AkianGaubertGuterman2012}.
There are also strongly polynomial claims based on tropical optimization
and max-atoms~\cite{Truffet2025MaxAtoms,Truffet2026SubstitutionV4}.
Appendix~\ref{sec:tropical-comparison} examines the printed substitution
rule in Truffet's arXiv version~4 of August~29, 2026 and gives an explicit
instance on which it loses the optimum. The exact theorem and procedure
text of the distinct 2025 HAL/MSR claim was unavailable to us; the
appendix makes no assertion about that proof. Its objection is specific
to the cited version and rule. Neither claim is an input to our proof.

\subsection{How the solver finds the potential}

First transform each weight to $(n+1)w(e)+1$. This familiar integer
perturbation separates the signs of simple-cycle sums
\cite[Section~9]{BjorklundVorobyov2007}. Given a vector $p$, let $F_i(p)$
be the maximum, at a Max vertex, or minimum, at a Min vertex, of a
transformed outgoing-edge weight plus the coordinate at its head.
The operator preserves order and commutes with common scalar shifts.
In a finite integer box $[a,b]=\{p\in\Z^V:a\le p\le b\}$,
with coordinatewise inequalities, impose $y_i\le F_i(y)$ only where
$y_i>a_i$ for a subsolution, and impose the reversed inequality only
where a supersolution lies below $b_i$. Every subsolution lies below
every supersolution. Clipping, or truncating each coordinate of $F$ to its interval, therefore
yields a unique
fixed point. In a sufficiently wide constant box its coordinates fall
into two separated bands, from which positional winning strategies
follow.

Ordinary monotone iteration can take time proportional to the box width.
We replace that dependence by recursion on the supports of the
comparisons. The support consists of the coordinates where the defining
inequality is imposed. Each side receives positive rational vertex
weights and its own budget. The recursive answer is constrained by its
order relative to each fixed admissible comparison; the answer itself
need not satisfy the comparison inequalities. Reducing a budget, or
halving weights on a suitable concentrated part of a support, decreases
a scalar rank parameter.

At a large box the solver first tightens its bounds, then chooses which
boundary to truncate. The auxiliary procedure called orientation returns
a choice of side and a list of vertex sets of small auxiliary measure.
For each fixed comparison whose support is sufficiently concentrated
near the threatened boundary, it either protects that side or returns a
set containing more than the required share of its support weight.
Reweighting on such a set makes a smaller-rank call applicable.

The orientation proof follows a sequence of centers. At each center,
two smaller-rank calls reduce the two budgets separately, and the next
center is their clipped average. Successful comparisons ensure that the
maximum deviation of the fixed comparison never grows. Whenever its
projection fits the reduced budget, that deviation decreases by a
controlled amount. Many such decreases force a visible displacement of
the extreme coordinates. If there are few, a newly sampled uniform
iterate is likely to have a heavy projected support in a smaller box.
The sampled comparison is fixed before the child's fresh random bits,
and support containment transports the child's conclusion back.

This order of conditioning is the reason pointwise guarantees suffice.
No event protecting all admissible vectors in one run is used.
Amplification controls the fresh smaller-rank calls. The one continuation
with unchanged support parameters is kept unamplified, and its entire
chain is analyzed together, both for errors and for work. At the top
box, the unique deterministic solution is an admissible comparison on
both sides. Intersecting its two whole-vector success events gives
exact recovery with probability at least $7/8$.

\subsection{Organization}

Section~\ref{sec:boxes} establishes finite-box comparison and the game
reduction. Section~\ref{sec:budgets} proves the support and projection
rules. Section~\ref{sec:algorithm} specifies the complete recursion and
proves termination on every tape. Under a lower-rank solver hypothesis,
Section~\ref{sec:orientation} proves orientation; Section~\ref{sec:solving}
then closes that rank induction and proves exact top-potential recovery.
Section~\ref{sec:complexity} bounds the full bit cost, and
Section~\ref{sec:completion} assembles the game algorithm and derives
the certifying and Las Vegas versions.
Appendix~\ref{sec:tropical-comparison} is an independent version-specific
comparison.
